%%%% ABSTRACT
%%
%% Versão do resumo para idioma de divulgação internacional.

\begin{abstractutfpr}%% Ambiente abstractutfpr

The demand for more efficient, cost-effective, and high-performance processors continues to grow, making research into architectural modifications aimed at improving energy efficiency and processing capability increasingly important. However, hardware design presents significant entry barriers due to its high complexity and associated costs. In this context, this work presents the planning and specification for the implementation of a processor core based on the RISC-V Instruction Set Architecture (ISA), specifically employing the 32-bit base integer instruction set (RV32I) in a single-cycle Harvard architecture. For the physical validation of the circuit, an AMD Artix-7 Field-Programmable Gate Array (FPGA), integrated into the Digilent Nexys A7 development board, will be adopted in order to establish a modular and educational platform that facilitates future investigations and application-specific customizations. To support the proposal, a literature review on processor development was conducted, identifying the growth of the RISC-V ecosystem across domains ranging from high-performance computing and Artificial Intelligence accelerators to low-power microcontrollers. Additionally, the structural concepts of the Harvard Architecture, its essential internal components, and the processing models of single-cycle, multi-cycle, and pipelined architectures were reviewed. The proposed methodology is systematically organized into seven consecutive stages, covering activities from architectural study and modular block modeling using the SystemVerilog hardware description language to final hardware testing. Finally, the partitioning of the system into five independent functional blocks and the established schedule provide the basis for project organization. As a result of this first stage of the undergraduate thesis, the work presents the initial planning and architectural mapping, contributing to the guidance and assessment of the technical feasibility of the future practical implementation and bench testing to be carried out in the next phase of the project.'

\end{abstractutfpr}
